\documentclass[11pt,a4paper]{article}

\usepackage[margin=2.5cm]{geometry}
\usepackage{amsmath}
\usepackage{xcolor}
\usepackage{listings}
\usepackage{booktabs}
\usepackage{enumitem}
\usepackage[colorlinks=true,urlcolor=blue,linkcolor=black]{hyperref}

\definecolor{codebg}{rgb}{0.96,0.96,0.96}
\definecolor{codekw}{rgb}{0.0,0.3,0.6}

\lstset{
  basicstyle=\ttfamily\small,
  backgroundcolor=\color{codebg},
  frame=single,
  framesep=4pt,
  breaklines=true,
  columns=fullflexible,
  keepspaces=true,
  showstringspaces=false,
}

\title{\textbf{Controlling a Single Dynamixel Servo}\\[4pt]
  \large A hands-on demo: raw DynamixelSDK and ROS~2}
\author{OMX\_jazzy workspace --- \texttt{dynamixel\_demo} package}
\date{\today}

\begin{document}
\maketitle

\begin{abstract}
\noindent This guide walks you through controlling a single ROBOTIS Dynamixel
X-series servo sitting on your bench. You will first drive it directly with the
\textbf{DynamixelSDK} (plain Python, no ROS), then with a minimal \textbf{ROS~2}
node. Both demos live in the \texttt{dynamixel\_demo} package of this workspace and
work with one servo wired to a U2D2 USB interface.
\end{abstract}

\tableofcontents
\vspace{1em}

% ---------------------------------------------------------------------------
\section{Overview}
The \texttt{dynamixel\_demo} package contains two self-contained demos:
\begin{description}[leftmargin=1.6cm,style=nextline]
  \item[Demo~1 --- \texttt{single\_servo\_demo.py}] Pure Python using the
    DynamixelSDK. It opens the serial port, \emph{pings} the servo to confirm
    wiring, enables torque, sweeps the horn between two positions while printing
    the measured position, then disables torque. No ROS required.
  \item[Demo~2 --- \texttt{servo\_ros\_node.py}] A minimal ROS~2 (\texttt{rclpy})
    node that opens the port itself and exposes two topics: it subscribes to a
    target angle on \texttt{/dxl/goal\_position} and publishes the measured angle
    on \texttt{/dxl/present\_position}. No MoveIt or ros2\_control needed.
\end{description}

Working through both shows the same hardware from two levels: the low-level byte
protocol, and the topic-based ROS~2 abstraction built on top of it.

% ---------------------------------------------------------------------------
\section{What you need (bill of materials)}
\begin{table}[h]
\centering
\begin{tabular}{@{}ll@{}}
\toprule
\textbf{Item} & \textbf{Notes} \\
\midrule
Dynamixel X-series servo & XL430-W250 or XM430-W350 (as used on the OMX). \\
U2D2 interface & USB-to-TTL/RS-485 adapter for Dynamixel. \\
U2D2 Power Hub Board & Or a 3-pin connection plus external power. \\
12\,V SMPS power supply & Powers the servo bus (X-series is 12\,V). \\
Robot cable (3-pin JST) & Connects the servo to the U2D2 / power hub. \\
USB cable & U2D2 to your computer. \\
\bottomrule
\end{tabular}
\caption{Hardware for the single-servo demo.}
\end{table}

\paragraph{Wiring.} Connect the 12\,V supply to the U2D2 Power Hub, daisy-chain
the servo to the hub with a robot cable, and plug the U2D2 into your computer over
USB. The servo has two identical 3-pin ports --- either may be used; the second is
for chaining further servos. \textbf{Always apply 12\,V power before expecting the
servo to respond}, and never hot-plug the power while torque is enabled.

% ---------------------------------------------------------------------------
\section{Finding the serial port and fixing permissions}
When you plug in the U2D2 it appears as a USB serial device, usually
\texttt{/dev/ttyUSB0}.
\begin{lstlisting}[language=bash]
ls /dev/ttyUSB*        # list candidate ports
dmesg | tail           # shows the device that was just attached
\end{lstlisting}

If running a demo prints a \emph{permission denied} error, add yourself to the
\texttt{dialout} group (then log out and back in):
\begin{lstlisting}[language=bash]
sudo usermod -aG dialout $USER
\end{lstlisting}

\paragraph{Baud rate and ID.} Every Dynamixel ships from the factory with
\textbf{ID~1} and a baud rate of \textbf{57600}. The demos default to these. If you
took the servo off an assembled Open Manipulator~X, its bus runs at
\textbf{1\,000\,000} baud instead --- pass \texttt{-{}-baud 1000000}. If a ping
fails, a wrong baud rate is the most common cause.

% ---------------------------------------------------------------------------
\section{Build and source the workspace}
\begin{lstlisting}[language=bash]
cd ~/Programming/OMX_jazzy
colcon build --packages-select dynamixel_demo --symlink-install \
    --allow-overriding dynamixel_sdk
source install/setup.bash
\end{lstlisting}
The \texttt{dynamixel\_sdk} Python module is already vendored in this workspace, so
no extra installation is required.

% ---------------------------------------------------------------------------
\section{Demo 1: direct control with the DynamixelSDK}
Run the standalone demo (no ROS core needed):
\begin{lstlisting}[language=bash]
ros2 run dynamixel_demo single_servo_demo --port /dev/ttyUSB0 --id 1 --baud 57600
# or simply:
python3 src/dynamixel_demo/dynamixel_demo/single_servo_demo.py --port /dev/ttyUSB0
\end{lstlisting}

Useful flags: \texttt{-{}-cycles N} (number of moves), \texttt{-{}-min} /
\texttt{-{}-max} (raw goal positions, $0\ldots4095$). Expected output:
\begin{lstlisting}
[ OK ] Opened port /dev/ttyUSB0
[ OK ] Baud rate set to 57600
[ OK ] Found Dynamixel ID 1 (model number 1060)
[ OK ] Torque enabled -- the servo is now holding position

  [ID:001] cycle 1/4  GoalPos:1024  PresPos:2047
  [ID:001] cycle 1/4  GoalPos:1024  PresPos:1031
  [ID:001] cycle 2/4  GoalPos:3072  PresPos:1040
  ...
[ OK ] Sweep complete.
[ OK ] Torque disabled -- the horn now turns freely.
[ OK ] Port closed.
\end{lstlisting}
The servo physically sweeps back and forth. When the demo ends it disables torque,
so the horn turns freely again.

% ---------------------------------------------------------------------------
\section{Demo 2: basic control with ROS 2}
Launch the node:
\begin{lstlisting}[language=bash]
ros2 launch dynamixel_demo servo_demo.launch.py port:=/dev/ttyUSB0 baud:=57600 dxl_id:=1
\end{lstlisting}

In a second terminal, command a target \textbf{angle in radians} (0 is the encoder
centre; the usable range is about $\pm\pi$):
\begin{lstlisting}[language=bash]
ros2 topic pub --once /dxl/goal_position std_msgs/msg/Float64 "{data: 0.5}"
ros2 topic pub --once /dxl/goal_position std_msgs/msg/Float64 "{data: -0.5}"
\end{lstlisting}

In a third terminal, watch the measured angle stream back:
\begin{lstlisting}[language=bash]
ros2 topic echo /dxl/present_position
\end{lstlisting}

\begin{table}[h]
\centering
\begin{tabular}{@{}lll@{}}
\toprule
\textbf{Topic} & \textbf{Type} & \textbf{Meaning} \\
\midrule
\texttt{/dxl/goal\_position} & \texttt{std\_msgs/Float64} & Target angle [rad] (you publish) \\
\texttt{/dxl/present\_position} & \texttt{std\_msgs/Float64} & Measured angle [rad] (node publishes) \\
\bottomrule
\end{tabular}
\caption{Topics exposed by the ROS~2 demo node.}
\end{table}
Stop the node with \texttt{Ctrl-C}; it disables torque and closes the port on exit.

% ---------------------------------------------------------------------------
\section{Troubleshooting}
\begin{table}[h]
\centering
\begin{tabular}{@{}p{5cm}p{8.5cm}@{}}
\toprule
\textbf{Symptom} & \textbf{Fix} \\
\midrule
\emph{Could not open port} & Wrong device name --- check \texttt{ls /dev/ttyUSB*}; pass the right \texttt{-{}-port}. \\
\emph{Permission denied} & Add yourself to the \texttt{dialout} group and re-login. \\
\emph{No response from ID~1} (ping fails) & Wrong baud (try \texttt{-{}-baud 1000000}), wrong ID, no 12\,V power, or a loose cable. \\
Servo does not move but ping works & Goal outside the servo's position limit, or it is in an error state --- power-cycle it. \\
Horn is limp / not holding & Torque is disabled (normal after a demo exits). \\
\bottomrule
\end{tabular}
\caption{Common problems and fixes.}
\end{table}

% ---------------------------------------------------------------------------
\appendix
\section{X-series control table reference}
The demos use the Protocol~2.0 control table shared by the XL430 and XM430:
\begin{table}[h]
\centering
\begin{tabular}{@{}llll@{}}
\toprule
\textbf{Item} & \textbf{Address} & \textbf{Size (bytes)} & \textbf{Notes} \\
\midrule
Torque Enable     & 64  & 1 & 1 = on, 0 = off \\
Goal Position     & 116 & 4 & Target, raw $0\ldots4095$ \\
Present Position  & 132 & 4 & Measured, raw $0\ldots4095$ \\
\bottomrule
\end{tabular}
\caption{Control-table addresses used by the demos.}
\end{table}
A full revolution spans 4096 raw units (a 12-bit encoder), centred at 2048. The
ROS~2 node converts between raw units and radians with
$\text{rad} = (\text{tick}-2048)\,/\,4096 \times 2\pi$.

\end{document}
